% ============================================================================
% chapters/ch2_background.tex
% ============================================================================
\chapter{Background and Related Work}
\label{ch:background}

This chapter reviews the theoretical foundations and prior work across four
domains that underpin this thesis: human activity recognition with deep
learning (\Cref{sec:bg_har}), MLOps practices for production machine learning
systems (\Cref{sec:bg_mlops}), domain adaptation and transfer learning for
sensor data (\Cref{sec:bg_domain_adaptation}), and uncertainty quantification
in deep neural classifiers (\Cref{sec:bg_uq}).

% ----------------------------------------------------------------------------
\section{Human Activity Recognition}
\label{sec:bg_har}
% ----------------------------------------------------------------------------

\subsection{Definition and Sensor Modalities}

Human Activity Recognition refers to the automatic identification of
physical activities performed by a person, inferred from sensor measurements
\citep{chen2021deep}.
In the wearable sensor paradigm relevant to this thesis, activities are
inferred from triaxial accelerometers and triaxial gyroscopes embedded in
a wrist-worn device.
The accelerometer measures linear acceleration along three orthogonal axes
(units: m/s\textsuperscript{2}), while the gyroscope measures angular velocity
around three orthogonal axes (units: $^\circ$/s or rad/s).

\citet{mannini2010machine} provide a foundational review of machine learning
methods for accelerometer-based HAR, establishing sensor validation ranges
that are still used in practice: acceleration $\leq 50$\,m/s\textsuperscript{2}
and gyroscope $\leq 500$\,$^\circ$/s.
\citet{bao2004activity} introduced the sliding window paradigm, observing that
a 4-second window captures sufficient context for coarse activity labels, a
convention adopted in most subsequent work.

\subsection{Deep Learning Architectures for HAR}

\citet{chen2021deep} survey deep learning approaches for sensor-based HAR,
identifying one-dimensional convolutional neural networks (1D-CNNs) and
recurrent architectures (LSTM, GRU) as the dominant families.
1D-CNNs extract local temporal features through learned filters, while
bidirectional LSTM layers capture global sequence context in both forward and
backward directions.
Their combination (\cnnbilstm{}) has been shown to outperform unimodal
architectures on benchmark datasets including UCI-HAR, PAMAP2, and
Opportunity.

The model used in this thesis is a \cnnbilstm{} with 499,131 parameters and
17 layers, pre-trained on labelled wrist sensor data using the 11 anxiety-related
activity classes described in \Cref{sec:scope}.
Details of the architecture are provided in \Cref{sec:model_architecture}.

\subsection{Benchmark Datasets and Limitations}

Most published HAR evaluations use subject-independent (leave-one-subject-out)
cross-validation on publicly available benchmarks.
A key limitation of this evaluation protocol is that it treats all subjects as
drawn from a single distribution, ignoring the persistent inter-subject
variability that arises in real deployments.
\todo{Add citations for benchmark datasets: UCI-HAR, PAMAP2, Opportunity.}

% ----------------------------------------------------------------------------
\section{MLOps for Production Machine Learning}
\label{sec:bg_mlops}
% ----------------------------------------------------------------------------

\subsection{The MLOps Problem}

\citet{sculley2015hidden} identified hidden technical debt in production ML
systems, observing that data dependencies, feedback loops, and system
entanglement generate maintenance costs that dwarf the cost of writing model
code.
Their analysis motivates the MLOps discipline: applying software engineering
principles (modularity, testing, monitoring, deployment automation) to the full
ML lifecycle, not just the modelling phase.

A canonical MLOps lifecycle comprises six phases: data management, model
development, model evaluation, deployment, monitoring, and retraining.
A production system must execute these phases repeatedly and automatically
as new data arrive \citep{sculley2015hidden}.

\subsection{Monitoring in Production}

The central challenge of production ML monitoring is detecting model
degradation without access to ground-truth labels.
\citet{gama2014survey} survey concept drift detection methods, identifying
three classes: sequential analysis methods (CUSUM, SPRT), distribution
comparison methods (KS test, PSI), and model-based methods (holdout accuracy
on labelled windows).
Only the first two classes are applicable in the unlabelled production setting
addressed by this thesis.

The Population Stability Index (PSI) is defined as:
\begin{equation}
  \text{PSI} = \sum_{i=1}^{B}
    \left( p_i^{\text{prod}} - p_i^{\text{ref}} \right)
    \ln \frac{p_i^{\text{prod}}}{p_i^{\text{ref}}}
  \label{eq:psi}
\end{equation}
where $p_i^{\text{prod}}$ and $p_i^{\text{ref}}$ are the production and
reference histogram proportions in bin $i$, and $B$ is the number of bins.
Standard thresholds ($\text{PSI} < 0.10$ no drift, $0.10$--$0.25$ moderate,
$> 0.25$ significant) apply to single-channel distributions; this thesis
empirically calibrates equivalent thresholds for 24-channel aggregated PSI
(see \Cref{sec:psi_calibration}).

\subsection{Data and Model Versioning}

\citet{sculley2015hidden} emphasise that data versioning is as important as
code versioning for reproducible ML.
DVC \citep{dvcteam2020} extends Git with pointer files that track large binary
artefacts (datasets and models) using content hashes, enabling reproducible
re-runs at any prior data version.
MLflow \citep{mlflowteam2018} provides a complementary layer: it records
experiment metadata, hyperparameters, metrics, and model artefacts associated
with code commits, linking DVC data versions to specific model versions.

% ----------------------------------------------------------------------------
\section{Domain Adaptation for Wearable Sensor Data}
\label{sec:bg_domain_adaptation}
% ----------------------------------------------------------------------------

\subsection{Sources of Domain Shift in HAR}

Domain shift in wearable HAR arises from three sources:
\begin{enumerate}
  \item \textbf{Inter-subject variability.} Different users move differently
        even when performing the same activity. A model trained on one user's
        data may not generalise to another user's distribution.
  \item \textbf{Sensor placement variability.} Wrist-worn sensors are affected
        by handedness (dominant vs.\ non-dominant wrist), tightness of fit,
        and wrist orientation relative to the device face.
  \item \textbf{Sensor calibration drift.} Long-term drift in accelerometer
        and gyroscope calibration changes the marginal distribution of sensor
        readings, even for the same physical activity.
\end{enumerate}

This thesis encountered a concrete instance of category~(3): production data
from Garmin smartwatches was recorded in milliG, while the pre-trained model
was trained on data in m/s\textsuperscript{2}.
The resulting magnitude difference degraded cross-user accuracy to 14.5\,\%
(below random-chance baseline of 9.09\,\% for 11 classes), a case study
detailed in \Cref{sec:unit_mismatch}.

\subsection{Adaptive Batch Normalisation}

\citet{li2018adabn} proposed Adaptive Batch Normalisation (AdaBN) as a
parameter-free test-time adaptation method.
The core observation is that Batch Normalisation \citep{ioffe2015batch} layers
store running statistics (mean and variance) that represent the training domain
distribution.
By replacing these statistics with the statistics computed from the target
domain, the model's internal normalisation is updated to reflect the new domain
without modifying any learnable weights.
AdaBN requires $n_{\text{batches}} = 10$ forward passes over the target data
to stabilise the running statistics, using batch size $B = 64$.

\subsection{TENT: Test-Time Entropy Minimisation}

\citet{wang2021tent} proposed TENT, a test-time adaptation method that
minimises the entropy of the model's output predictions over a target batch
through gradient descent on the Batch Normalisation affine parameters
($\gamma, \beta$) only.
TENT uses $n_{\text{steps}} = 10$ gradient steps with learning rate
$lr = 10^{-4}$.

An important interaction arises when AdaBN and TENT are applied sequentially:
TENT's training-mode forward passes update the BN running statistics, overwriting
the statistics set by AdaBN.
\Cref{sec:tent_adabn} describes the snapshot-and-restore mechanism implemented
in this thesis to prevent this interaction.

\subsection{Pseudo-Labelling for Self-Supervised Adaptation}

\citet{lee2013pseudo} introduced pseudo-labelling as a semi-supervised learning
strategy: a model trained on labelled data generates pseudo-labels for unlabelled
examples; those pseudo-labels are then used as training targets.
Quality is controlled by a confidence threshold $\tau$, retaining only
predictions with maximum softmax probability $> \tau$.
The curriculum variant progressively lowers $\tau$ across training iterations,
providing easy-to-classify examples first \citep{lee2013pseudo}.

Catastrophic forgetting is a risk in any fine-tuning scenario: the model may
overfit to the new domain and lose performance on the original domain.
\citet{kirkpatrick2017overcoming} addressed this with Elastic Weight
Consolidation (EWC), which adds an $L_2$ penalty proportional to the Fisher
information to anchor weights that were important for the original task.

% ----------------------------------------------------------------------------
\section{Uncertainty Quantification}
\label{sec:bg_uq}
% ----------------------------------------------------------------------------

\subsection{Calibration and Expected Calibration Error}

A neural network is said to be \textit{calibrated} if its predicted confidence
values correspond to empirical accuracy: a batch of predictions with 80\,\%
confidence should be correct roughly 80\,\% of the time.
Modern deep neural networks are typically overconfident \citep{guo2017calibration}.

\citet{guo2017calibration} showed that temperature scaling --- dividing the
pre-softmax logit vector by a scalar temperature $T > 1$ --- effectively
recalibrates neural networks with minimal additional computation.
The optimal temperature is found by minimising Expected Calibration Error (ECE)
on a held-out validation set.
ECE is computed as:
\begin{equation}
  \text{ECE} = \sum_{b=1}^{B} \frac{|S_b|}{N}
    \left| \mathrm{acc}(S_b) - \mathrm{conf}(S_b) \right|
  \label{eq:ece}
\end{equation}
where $S_b$ is the set of samples in confidence bin $b$, $|S_b|/N$ is the
fraction of samples in that bin, and $\mathrm{acc}$, $\mathrm{conf}$ are the
bin accuracy and mean confidence, respectively.
A warning is raised when $\text{ECE} > 0.10$ \citep{naeini2015obtaining}.

\subsection{Monte Carlo Dropout}

\citet{gal2016dropout} showed that applying dropout at test time (keeping
dropout active during inference) approximates Bayesian inference in a deep
Gaussian process.
Running $T$ stochastic forward passes through the model with dropout enabled
produces a distribution over predictions; the variance of this distribution
is an estimator of epistemic (model) uncertainty.
$T = 30$ passes are recommended for stable uncertainty estimates
\citep{gal2016dropout}.
